\scanpage{66}
\[
\unsure{\mathcal X_0|\mathcal X|\mathfrak X_0}\longmapsto\underline{\operatorname{Isom}}(\mathcal X_0,\mathcal X)\tag*{(15)}
\]
vers \deleted{les} $\operatorname{Tors}(\underline{\operatorname{Aut}}(\mathcal X_{0S}))$, avec l'équivalence
\[
\operatorname{Tors}(\underline{\operatorname{Aut}}(\mathcal X_{0S}))\simeq\operatorname{Tors}(\mathrm{GP}(1)_S)\tag*{(16)}
\]
déduite de l'iso \deleted{\unsure{(7)}}
\[
\mathrm{GP}(1)_S\xrightarrow{\sim}\underline{\operatorname{Aut}}(\mathcal X_{0S})\bigl(\simeq\underline{\operatorname{Aut}}(\mathcal X_0)_S\bigr)\tag*{(17)}
\]
déduit de (7) par chgt de base.

\inserted{\unsure{Également} forme \unsure{particulière} du foncteur ``twist'' dans le cas de $C_{0S}\to C_{1S}$, cf ci-dessous}

\uline{IV} Le foncteur allant de la catégorie \deleted{$C_{1S}$ des formes de $\mathrm{GP}(1)_S$ vers des catégories} \deleted{catégories} $C_{0S}, C_{2S}$ etc vers la catégorie \inserted{$C_{1S}$} des formes de $\mathrm{GP}(1)$ s'explicite par
\[
\mathcal X\longrightarrow\underline{\operatorname{Aut}}(\mathcal X)\tag*{(18)}
\]

\deleted{Pour expliciter les foncteurs} Dans le cas \unsure{de|des} $C_{1S}'$ des formes $\widetilde G$ de $\mathrm{SL}(2)$, il y a aussi une description plus terre à terre
\[
\widetilde G\longrightarrow\widetilde G/\operatorname{Cent}(\widetilde G)\tag*{(19)}
\]
(\uline{NB} le centre de $\widetilde G$ est tjrs \deleted{\unsure{can}} isom à $\mu_2$ (de façon unique car $\operatorname{Aut}(\mu_2)=\overline1$

Pour le cas des $\underline{\operatorname{Quat}}(\mathcal O_S)$ $\to\operatorname{For}(G)\ldots)$
\[
M\longmapsto\check{\mathbb V}(M)^*/G_{mS}\tag*{(20)}
\]
% archive sheet number: 65
